\section{从 ReAct 到 Plan-and-Solve}

\subsection{ReAct 的局限性}

ReAct 通过 Thought $\rightarrow$ Action $\rightarrow$ Observation 循环来处理复杂问题，但存在一个关键问题：

\begin{warningbox}{Context Issues}
当 ReAct 循环特别长时，会出现严重的上下文问题：
\begin{itemize}
    \item 语言模型会 \textbf{Lost in the Middle}——迷失在冗长的上下文中
    \item 忘记用户的原始任务是什么
    \item 忘记自己做过哪些事情、得到了什么结果
    \item 无法产生下一步合理的 Action
\end{itemize}
\end{warningbox}

\subsection{Plan-and-Solve Prompting}

Plan-and-Solve 是对经典 Zero-shot CoT 的改进。其 Prompting 方式为：

\begin{quote}
\textit{``Let's first understand the problem and devise a step-by-step plan to solve it. Then let's carry out the plan and solve the problem step by step.''}
\end{quote}

对比经典的 ``Let's think step by step''，Plan-and-Solve 显式地在 Prompting 中引入了 Plan 的概念，在复杂问题上大约有几个百分点的提升。

\subsection{Dynamic Plan-and-Solve}

社区对 Plan-and-Solve 做了进一步改进，引入了\textbf{动态性}：

\begin{importantbox}{刚中之脑 vs. 真实交互}
Plan 是在头脑中形成的（刚中之脑），没有和环境真实交互。当真实执行时会产生新的 feedback，这时需要 \textbf{Re-Plan}——将新的反馈引入到规划机制中进行动态修正。
\end{importantbox}

Dynamic Plan-and-Solve 的完整流程：
\begin{enumerate}
    \item \textbf{Plan}：Planner 自顶向下形成一个 plan（List of Steps / To-Do List）
    \item \textbf{Execute}：每个 step 交给一个 ReAct Agent 去执行
    \item \textbf{Re-Plan}：执行后基于新的 feedback，决定是直接结束还是重新规划
\end{enumerate}

Re-Plan 的输入包括：原始问题 + 原始 Plan + 过去执行的步骤及其结果。

\begin{knowledgebox}{经济性考量：强弱模型组合}
从经济角度，可以用\textbf{强模型}做 Planner 和 Re-Plan（因为规划决定了整体质量上限），用\textbf{弱模型}做具体的 ReAct 执行（因为具体执行步骤相对简单）。这样在成本和效果之间取得平衡。
\end{knowledgebox}

\subsection{本章小结}

从 ReAct 到 Plan-and-Solve 再到 Dynamic Plan-and-Solve，体现了一条清晰的演进线：从纯反应式交互，到显式引入规划，再到动态调整规划以弥合"刚中之脑"与真实交互之间的 gap。


